% SGA 3，第 VII_A 讲， 英文重建稿——第 1 轮。
% 范围：权威底本局部页 40--44, 全部 节 6,
% 止于其前： 节 7 局部页 45.
% 合并阅读本页码： 493--497. 再版印刷页码： 482--486.
% 本范围内原版页边定位： 454--460.
% 权威底本： Polo--Gille Exp7A-13oct24.pdf,
% SHA-256 D46DB8C79D0385643B6EED4B96900613248375F9B73B4202478A4F57FD3E8522.
% 曾参考 Jacob Reinhold 的 CC BY 4.0 Markdown 作为比对谱系。

\section*{6. \(S\)-群概形的 \(p\)-Lie 代数}
\phantomsection
\label{sga3:VIIA:6}
\addcontentsline{toc}{section}{6. S-群概形的 p-Lie 代数}

% 权威底本： Exp7A 局部页 40; 合并本页 493; 印刷页 482.
设 \(S\) 是特征为 \(p>0\) 的概形。在 \SGARefLink{sga3:VIIA:5.5}{5.5} 中，我们给每个拟凝聚 \(\mathcal O_S\)-\(p\)-Lie 代数 \(\mathcal L\) 配备了群 \(S\)-函子
\[
  \mathfrak G_p(\mathcal L)
  =
  \operatorname{Spec}^{*}\mathcal U_p(\mathcal L).
\]
现在将看到，对于每个 \(S\)-群概形 \(G\)，II，\SGARefLink{sga3:ii:subsection:4-11}{4.11} 中定义的 \(\mathcal O_S\)-Lie 代数 \(\operatorname{Lie}(G/S)\) 自然地配有 \(\mathcal O_S\)-\(p\)-Lie 代数结构。

\medskip
\noindent\textbf{6.1.}
\phantomsection
\label{sga3:VIIA:6.1}

首先借助 \SGARefLink{sga3:VIIA:2.5-proposition}{2.5} 的单射 \(\alpha\) 和 \(\beta\)，分别将 \(\operatorname{Lie}(G/S)(S)\) 和 \(\operatorname{Lie}(\operatorname{Aut}G/S)(S)\) 等同于 \(\mathrm U(G)\) 和 \(\operatorname{Dif}_{G/S}\) 的 Lie 子代数。于是 \(\operatorname{Lie}(\operatorname{Aut}G/S)(S)\) 等同于 \(\mathcal O_G\) 的 \(S\)-导子所成的 \(\Gamma(S,\mathcal O_S)\)-Lie 代数。由 \SGARefLink{sga3:VIIA:5.2}{5.2}，后者是 \(\operatorname{Dif}_{G/S}\) 的受限 \(p\)-Lie 子代数。

另一方面，
\[
  L=\operatorname{Lie}(G/S)(S)
\]
在单的代数同态
\[
  \ell:\mathrm U(G)\longrightarrow\operatorname{Dif}_{G/S},
  \qquad
  d\longmapsto {}_Gd,
\]
下的像由左不变导子组成（参见 \SGARefLink{sga3:VIIA:2.2}{2.2}、编者注~(17)、\SGARefLink{sga3:VIIA:2.4}{2.4} 和 \SGARefLink{sga3:VIIA:2.5-proposition}{2.5}）。若 \(x\in L\)，则由所引段落有 \(\ell(x)^p=\ell(x^p)\)。由于 \(\ell(x)^p\) 仍是导子，\(x^p\) 属于 \(\operatorname{Lie}(G/S)(S)\)。因此
\begin{center}
  \textit{\(\operatorname{Lie}(G/S)(S)\) 是无穷小代数 \(\mathrm U(G)\) 的受限 \(p\)-Lie 子代数。}
\end{center}
\footnote{也可以不经过 \(\operatorname{Dif}_{G/S}\)，直接证明 \(G\) 在原点处的导子所成的 Lie 代数（由 \SGARefLink{sga3:VIIA:2.5-proposition}{2.5}，它同构于 \(\operatorname{Lie}(G/S)(S)\)）在 \(\mathrm U(G)\) 中取 \(p\) 次幂下稳定。下文 \SGARefLink{sga3:VIIA:6.2}{6.2} 将给出这一证明。}

\medskip
\noindent\textbf{6.1.1.}
\phantomsection
\label{sga3:VIIA:6.1.1}

% 权威底本： Exp7A 局部页 41; 合并本页 494; 印刷页 483.
设 \(\varphi:G\to H\) 是 \(S\)-群概形同态。同态 \(\operatorname{Lie}(\varphi/S)(S)\) 和 \(\mathrm U(\varphi)\) 与将 \(\operatorname{Lie}(G/S)(S)\) 和 \(\operatorname{Lie}(H/S)(S)\) 等同于 \(\mathrm U(G)\) 和 \(\mathrm U(H)\) 的受限 \(p\)-Lie 子代数的做法相容。由于 \(\mathrm U(\varphi)\) 是代数同态，\(\operatorname{Lie}(\varphi/S)(S)\) 是受限 \(p\)-Lie 代数同态。

同样，若 \(s:T\to S\) 是基变换，则由 \(s\) 诱导的从 \(\operatorname{Lie}(G/S)(S)\) 到 \(\operatorname{Lie}(G/S)(T)\) 的映射是受限 \(p\)-Lie 代数同态。等价地，函子 \(\operatorname{Lie}(G/S)\) 配有 \(\mathcal O_S\)-\(p\)-Lie 代数结构。特别地，当 \(T\) 遍历 \(S\) 的开子集时，层 \(\operatorname{Lie}(G/S)\) 配有 \(\mathcal O_S\)-\(p\)-Lie 代数结构。

\medskip
\noindent\textbf{6.2.}
\phantomsection
\label{sga3:VIIA:6.2}

按照 Demazure 的一个想法，我们将把上述构造推广到某些未必可表的群 \(S\)-函子。首先给出 \(S\)-群概形 \(G\) 的 Lie 代数中符号 \(p\) 次幂的另一个定义。

设 \(D\) 是 \(G\) 在原点处的导子。\footnote{在下文中，编者扩充了原文的论证。}由 \SGARefLink{sga3:VIIA:1.2.1}{1.2.1}，\(D\) 是零截面 \(\tau:S\to I_S\) 的 \(S\)-导子
\[
  \delta=(\tau,\partial_t)
\]
与一个满足 \(x\circ\tau=\varepsilon\) 的态射 \(x:I_S\to G\) 的复合，也就是说，\(x\in\operatorname{Lie}(G/S)(S)\)。由 \SGARefLink{sga3:VIIA:2.1}{2.1} 中的定义，\(D^p\) 是如下复合偏差：
\begin{center}
  \phantomsection
  \label{sga3:VIIA:6.2-demazure-pfold-composite}
  % SGA3-VII-DIAG-057
  \resizebox{0.98\linewidth}{!}{%
  \(\displaystyle
  S \simeq
  \underbrace{S\times S\times\cdots\times S}_{p}
    \xrightarrow{\delta\times\cdots\times\delta}
  I_S\times\cdots\times I_S
    \xrightarrow{x\times\cdots\times x}
  G\times\cdots\times G
    \xrightarrow{m^{(p)}} G
  \)%
}

\end{center}
这里 \(m^{(p)}\) 由乘法 \(m:G\times G\to G\) 诱导。由于 \(I_S\times\cdots\times I_S\) 在 \(S\) 上是仿射的，其仿射代数为
\[
  \mathcal B
  =
  \mathcal O_S[d_1,\ldots,d_p]/
  (d_1^2,\ldots,d_p^2),
\]
，偏差 \(\delta\times\cdots\times\delta\) 由如下 \(\mathcal O_S\)-模同态定义：
\[
  \phi:\mathcal B\longrightarrow\mathcal O_S
\]
它将单项式 \(d_1d_2\cdots d_p\) 映至 \(1\)，而将所有其他单项式 \(d_{i_1}\cdots d_{i_r}\)（\(0\le r<p\)）映至 \(0\)。

以 \(\operatorname{pr}_i\) 表示 \(I_S^p\) 到其第 \(i\) 个因子的投影，并令 \(x_i\) 为 \(x\) 在 \(G(\operatorname{pr}_i)\) 下在 \(G(I_S^p)\) 中的像。则
\[
  m^{(p)}\circ(x\times\cdots\times x)
\]
恰好是乘积 \(x_1x_2\cdots x_p\)。因而，\(D^p\) 也是复合偏差
\begin{center}
  \phantomsection
  \label{sga3:VIIA:6.2-demazure-product-composite}
  % SGA3-VII-DIAG-058
  \(\displaystyle
S \xrightarrow{\delta\times\cdots\times\delta}
I_S\times\cdots\times I_S
\xrightarrow{x_1x_2\cdots x_p}
G .
\)

\end{center}

这一描述给出了 \(D^p\) 是 \(G\) 在原点处导子的另一个证明。由于 \(G\) 是甚好群（II，\SGARefLink{sga3:ii:subsection:4-11}{4.11}），\(x\) 在 \(G(I_S\times I_S)\) 中的像 \(G(\operatorname{pr}_1)(x)\) 和 \(G(\operatorname{pr}_2)(x)\) 可交换。由此可知，元素 \(x_i\in G(I_S^p)\) 两两可交换。因此，对于 \(I_S^p\) 的因子的每个置换 \(\gamma\)，有
\[
  (x_1\cdots x_p)\circ\gamma=x_1\cdots x_p.
\]
因此 \(x_1\cdots x_p\) 通过从 \(I_S^p\) 到对称积 \(\Sigma^p I_S\) 的典范投影分解（参见 \SGARefLink{sga3:VIIA:4.2}{4.2}）。

\(\Sigma^p I_S\) 的仿射代数是 \(\mathcal B\) 的子代数 \(\mathcal A\)，它以 \(d_1,\ldots,d_p\) 的初等对称函数
\[
  1=\sigma_0,\sigma_1,\ldots,\sigma_p
\]
为 \(\mathcal O_S\)-基。令 \(\kappa:\mathcal A\hookrightarrow\mathcal B\) 为包含映射，并令
\[
  \pi:\mathcal A\longrightarrow\mathcal O_S[t]/(t^2)
\]
为将 \(\sigma_1,\ldots,\sigma_{p-1}\) 映为零、将 \(\sigma_p\) 映为 \(t\) 的 \(\mathcal O_S\)-代数同态。则
\[
  \phi\circ\kappa=\partial_t\circ\pi,
\]
其中 \(\partial_t:\mathcal O_S[t]/(t^2)\to\mathcal O_S\) 将 \(1\) 映为零，将 \(t\) 映为 \(1\)。若
\[
  i:I_S\hookrightarrow\Sigma^p I_S
\]
表示由 \(\pi\) 定义的闭浸入，则得到交换图
\begin{center}
  \phantomsection
  \label{sga3:VIIA:6.2-symbolic-power-commutative-diagram}
  % SGA3-VII-DIAG-059
  \begin{tikzcd}[row sep=3.5em,column sep=4.3em]
S
  \arrow[r,swap,"\delta\times\cdots\times\delta"]
  \arrow[d,swap,"\delta"]
  \arrow[rr,bend left=23,"D^p"]
  & I_S^p \arrow[r,swap,"x_1\cdots x_p"] \arrow[d,"\text{典范}"]
  & G \arrow[d,equal] \\
I_S
  \arrow[r,"i"]
  \arrow[rr,bend right=23,swap,"y"]
  & \Sigma^p I_S \arrow[r]
  & G
\end{tikzcd}

\end{center}
，这表明 \(D^p\) 具有 \(y\circ\delta\) 的形式，因而确实是 \(G\) 在原点处的导子。

\medskip
\noindent\textbf{6.3.}
\phantomsection
\label{sga3:VIIA:6.3}

% 权威底本： Exp7A 局部页 42; 合并本页 495; 印刷页 484.
令 \(\mathfrak S_p\) 为 \(p\) 次对称群，以 \(I_S^p\times\mathfrak S_p\) 表示一族以 \(\mathfrak S_p\) 为指标的 \(I_S^p\) 的副本的不交并。令
\[
  \pi:I_S^p\times\mathfrak S_p\longrightarrow I_S^p
\]
为典范投影，并令
\[
  \mu:I_S^p\times\mathfrak S_p\longrightarrow I_S^p
\]
为定义 \(\mathfrak S_p\) 在 \(I_S^p\) 上作用的态射：若 \(\rho\in\mathfrak S_p\)，则 \(\mu\) 在 \(I_S^p\times\rho\) 上的限制的第 \(j\) 个分量为 \(\operatorname{pr}_{\rho(j)}\)。

\noindent\textbf{定义。}
称函子
\[
  X:(\operatorname{Sch}/S)^\circ\longrightarrow(\operatorname{Ens})
\]
满足\emph{条件 (F)}，如果：
\begin{enumerate}
  \item[\textit{a)}] \(X\) 将有限不交并变为直积；
  \item[\textit{b)}] 对于每个 \(S\)-概形 \(T\)，下面的列正合：
\end{enumerate}
\begin{center}
  \phantomsection
  \label{sga3:VIIA:6.3-conditionF-equalizer}
  % SGA3-VII-DIAG-060
  \resizebox{0.98\linewidth}{!}{%
\begin{tikzcd}[ampersand replacement=\&,column sep=4.4em]
X(T\times\Sigma^p I_S)
  \arrow[r]
  \& X(T\times I_S^p)
    \arrow[r,shift left=.8ex,"X(\operatorname{id}_T\times\pi)"]
    \arrow[r,shift right=.8ex,swap,"X(\operatorname{id}_T\times\mu)"]
  \& X(T\times I_S^p\times\mathfrak S_p)\, .
\end{tikzcd}%
}

\end{center}

每个 \(S\)-概形都满足 (F)。若 \(\mathcal F\) 是 \(\mathcal O_S\)-模，则 \(\mathcal W(\mathcal F)\) 满足 (F)。满足 (F) 的函子的每个投射极限仍满足 (F)。最后，若 \(Y\) 满足 (F)，而 \(X\) 是任意 \(S\)-函子，则 \(\operatorname{Hom}_S(X,Y)\) 满足 (F)。

设 \(X\) 是满足条件 (F) 的甚好群（II，\SGARefLink{sga3:ii:definition:4-10}{4.10}）。设 \(x:I_S\to X\) 是延拓 \(X\) 的单位截面的态射，并采用 \SGARefLink{sga3:VIIA:6.2}{6.2} 的记号。与上面一样，
\[
  x_1\cdots x_p:I_S^p\longrightarrow X
\]
通过 \(\Sigma^p I_S\) 分解：
\begin{center}
  \phantomsection
  \label{sga3:VIIA:6.3-symbolic-power-factorization}
  % SGA3-VII-DIAG-061
  \begin{tikzcd}[row sep=3.5em,column sep=2.1em]
I_S^p \arrow[rr,"x_1\cdots x_p"] \arrow[dr,swap,"\text{典范}"]
  && X \\
& \Sigma^p I_S \arrow[ur,swap,"\Sigma^p(x)"]
&
\end{tikzcd}

\end{center}
因而通过复合定义了一个态射
\begin{center}
  \phantomsection
  \label{sga3:VIIA:6.3-symbolic-pth-power-composite}
  % SGA3-VII-DIAG-062
  \(\displaystyle
x^{(p)}:\quad
I_S \xrightarrow{i} \Sigma^p I_S
  \xrightarrow{\Sigma^p(x)} X
\)

\end{center}
，我们称之为 \(x\) 的\emph{符号 \(p\) 次幂}。

% 权威底本： Exp7A 局部页 43; 合并本页 496; 印刷页 485.
\(\operatorname{Lie}(X/S)(S)\) 的自同态
\[
  x\longmapsto x^{(p)}
\]
显然与基变换相容，并且关于 \(X\) 具有函子性。确定对于哪些 \(X\)，这一自同态使 \(\operatorname{Lie}(X/S)(S)\) 成为受限 \(p\)-Lie 代数，将是有意义的。

\medskip
\noindent\textbf{6.4.}
\phantomsection
\label{sga3:VIIA:6.4}

上述符号 \(p\) 次幂的定义特别适合计算。下面给出一些例子。

\medskip
\noindent\textbf{6.4.1.}
\phantomsection
\label{sga3:VIIA:6.4.1}

设 \(M\) 是一个“抽象”阿贝尔群，令 \(D_S(M)\) 为类型为 \(M\) 的可对角化 \(S\)-群（I，\SGARefLink{sga3:i:paragraph:4-4-2}{4.4.2}）。对于每个 \(S\)-概形 \(T\)，
\[
  D_S(M)(T)
  =
  \operatorname{Hom}_{(\mathrm{Ab})}
  \bigl(M,\mathcal O(T)^\times\bigr).
\]
设 \(x\in\operatorname{Lie}(D_S(M)/S)(S)\)，即一个阿贝尔群同态
\[
  x:M\longrightarrow
  \Gamma(S,\mathcal O_S+d\mathcal O_S)^\times
\]
，其形式为
\[
  m\longmapsto1+d\xi(m),
  \qquad
  \xi\in
  \operatorname{Hom}_{(\mathrm{Ab})}(M,\mathcal O(S)).
\]
沿用 \SGARefLink{sga3:VIIA:6.2}{6.2} 和 \SGARefLink{sga3:VIIA:6.3}{6.3} 的记号，乘积 \(x_1\cdots x_p\) 将 \(m\in M\) 对应到表达式
\[
  (1+d_1\xi(m))\cdots(1+d_p\xi(m)),
\]
，即
\[
  1+\sigma_1\xi(m)+\sigma_2\xi(m)^2+\cdots+
  \sigma_p\xi(m)^p.
\]
此表达式属于 \(\mathcal O(\Sigma^p I_S)\)。将其投影到 \(\mathcal O(S)[d]/(d^2)\)，令 \(\sigma_1,\ldots,\sigma_{p-1}\) 为零，并将 \(\sigma_p\) 映至 \(d\)，便可看出 \(x^{(p)}\) 是同态
\[
  m\longmapsto1+d\xi(m)^p.
\]
因此，在 II，\SGARefLink{sga3:ii:proposition:5-1}{5.1} 的等同
\[
  \operatorname{Lie}(D_S(M)/S)(S)
  \simeq
  \operatorname{Hom}_{(\mathrm{Ab})}(M,\mathcal O(S))
\]
下，符号 \(p\) 次幂将 \(\xi\) 映至
\[
  \xi^{(p)}:m\longmapsto\xi(m)^p.
\]

\medskip
\noindent\textbf{6.4.2.}
\phantomsection
\label{sga3:VIIA:6.4.2}

设 \(\mathcal F\) 是 \(\mathcal O_S\)-模，令 \(G=\mathcal W(\mathcal F)\) 为取值于阿贝尔群的群 \(S\)-函子（参见 I，\SGARefLink{sga3:i:subsection:4-6}{4.6}）。设
\[
  y\in\mathcal W(\mathcal F)(S)=\Gamma(S,\mathcal F),
\]
并令 \(y'\) 为它在 \(\mathcal W(\mathcal F)(I_S\to S)\) 下在 \(\mathcal W(\mathcal F)(I_S)\) 中的像。

由 II，\SGARefLink{sga3:ii:examples:4-4-2}{4.4.2} 和 \SGARefLink{sga3:ii:scholium:4-5-1}{4.5.1}，映射
\[
  y\longmapsto dy'
\]
是从 \(\mathcal W(\mathcal F)(S)\) 到 \(\operatorname{Lie}(\mathcal W(\mathcal F)/S)(S)\) 的 \(\mathcal O(S)\)-模同构。令 \(x=dy'\)。则 \SGARefLink{sga3:VIIA:6.2}{6.2} 中的元素 \(x_i\) 为 \(d_i y''\)，其中 \(y''\) 表示 \(y'\)\footnote{编者将 \(x\) 改为 \(y'\)。}在 \(\mathcal W(\mathcal F)(I_S^p)\) 中的典范像。因此
\[
  x_1+\cdots+x_p
  =
  (d_1+\cdots+d_p)y''
  =
  \sigma_1y'',
\]
且它属于 \(\mathcal W(\mathcal F)(\Sigma^p I_S)\)。由于定义态射 \(i\) 的同态
\[
  \mathcal O(\Sigma^p I_S)\longrightarrow\mathcal O(I_S)
\]
将 \(\sigma_1\) 映为零，有 \(x^{(p)}=0\)。因此
\begin{center}
  \textit{对于每个 \(\mathcal O_S\)-模 \(\mathcal F\)，\(\operatorname{Lie}\mathcal W(\mathcal F)\) 上的运算 \(x\mapsto x^{(p)}\) 为零。}
\end{center}

\medskip
\noindent\textbf{6.4.3.}
\phantomsection
\label{sga3:VIIA:6.4.3}

% 权威底本： Exp7A 局部页 44; 合并本页 497; 印刷页 486.
设 \(X\) 是 \(S\)-概形，令 \(G=\operatorname{Aut}_S X\)，并设 \(D\) 是结构层 \(\mathcal O_X\) 的 \(S\)-导子。由 \SGARefLink{sga3:VIIA:1.5}{1.5}，可将 \(D\) 等同于 \(X_{I_S}\) 的一个 \(I_S\)-自同构 \(x\)，它在 \(X\) 上诱导恒等映射。其描述如下。若 \(f\) 是 \(\mathcal O_X[d]/(d^2)\) 的一个形如 \(a+bd\) 的截面，令
\[
  D_{I_S}f=Da+(Db)d.
\]
于是 \(D_{I_S}\) 由 \(D\) 通过基变换 \(I_S\to S\) 得到，而 \(X_{I_S}\) 的这个自同构对应于 \(\mathcal O_X[d]/(d^2)\) 的自同态
\[
  f\longmapsto f+(D_{I_S}f)d
  =
  a+(D(a)+b)d
\]
。

同样，令 \(D_{I_S^p}\) 为由 \(D\) 通过基变换 \(I_S^p\to S\) 得到的 \(X_{I_S^p}\) 上的微分算子。沿用 \SGARefLink{sga3:VIIA:6.2}{6.2} 的记号，\(X_{I_S^p}\) 的自同构 \(x_i\) 对应于如下自同态
\[
  f\longmapsto f+(D_{I_S^p}f)d_i
\]
，其作用在
\[
  \mathcal O_X[d_1,\ldots,d_p]/(d_1^2,\ldots,d_p^2).
\]
上。因此乘积 \(x_1\cdots x_p\) 对应于
\[
  (1+d_1D_{I_S^p})(1+d_2D_{I_S^p})\cdots
  (1+d_pD_{I_S^p}),
\]
，即对应于
\[
  1+\sigma_1D_{I_S^p}
  +\sigma_2D_{I_S^p}^{\,2}
  +\cdots+
  \sigma_pD_{I_S^p}^{\,p}.
\]
\(\sigma_p\) 的系数为 \(D_{I_S^p}^{\,p}\)。因而，Lie 代数同构
\[
  \operatorname{Der}_S(\mathcal O_X)
  \xrightarrow{\ \sim\ }
  \operatorname{Lie}(\operatorname{Aut}_S X),
  \qquad
  D\longmapsto x
\]
（参见 \SGARefLink{sga3:VIIA:1.5}{1.5}）也是受限 \(p\)-Lie 代数的同构。

\medskip
\noindent\textbf{6.4.4.}
\phantomsection
\label{sga3:VIIA:6.4.4}

同一方法表明，对于每个 \(\mathcal O_S\)-模 \(\mathcal F\)，Lie 代数同构
\[
  \operatorname{Lie}
  \bigl(
    \operatorname{Aut}_{\mathcal O_S\text{-mod.}}
    \mathcal W(\mathcal F)/S
  \bigr)(S)
  \xrightarrow{\ \sim\ }
  \bigl(
    \operatorname{End}_{\mathcal O_S\text{-mod.}}
    \mathcal W(\mathcal F)
  \bigr)(S)
\]
（参见 II，\SGARefLink{sga3:ii:proposition:4-8}{4.8}）也是受限 \(p\)-Lie 代数的同构。

\medskip
\noindent\textbf{6.4.5.}\footnote{编者补上了编号 \SGARefLink{sga3:VIIA:6.4.5}{6.4.5}，并扩充了原文。}
\phantomsection
\label{sga3:VIIA:6.4.5}

设 \(\mathcal U\) 是群 \(\mathcal O_S\)-余代数，并令
\[
  G=\operatorname{Spec}^{*}\mathcal U
\]
为相应的群 \(S\)-函子；假设 \(G\) 可表。对于每个 \(T\to S\)，\SGARefLink{sga3:VIIA:5.5}{5.5} 和 \SGARefLink{sga3:VIIA:6.1.1}{6.1.1} 在
\[
  L(T)=\operatorname{Lie}(G)(T).
\]
上定义了两个受限 \(p\)-Lie 代数结构。存在交换图
\begin{center}
  \phantomsection
  \label{sga3:VIIA:6.4.5-group-coalgebra-comparison}
  % SGA3-VII-DIAG-063
  \begin{tikzcd}[row sep=3.4em,column sep=4.3em]
L(T) \arrow[r,hook,"\tau"] \arrow[d,hook,swap,"\alpha"] \arrow[dr,"i"]
  & \Gamma(T,\mathcal U_T) \\
\mathrm U(G_T)
  & \mathrm U(L(T)) \arrow[l,swap,"\varphi"] \arrow[u,swap,"\psi"]
\end{tikzcd}

\end{center}
其中 \(\mathrm U(L(T))\) 是 \(L(T)\) 的包络代数，而 \(\varphi,\psi\) 是由 \(\alpha,\tau\) 诱导的代数同态。由此可知，这两个受限 \(p\)-Lie 代数结构相同：若将 \(x\in L(T)\) 分别等同于它在 \(\mathrm U(G_T)\) 和 \(\Gamma(T,\mathcal U_T)\) 中的像，则 \(x^{(p)}\) 分别是元素 \(x^p\in\mathrm U(L(T))\) 在 \(\varphi\) 和 \(\psi\) 下的像。

% 下一游标： Exp7A 局部页 45, 合并阅读本页 498,
% 再版印刷页 487, 位于 节 7.
